%----------------------------------------------------------------------------------------
%	PACKAGES AND OTHER DOCUMENT CONFIGURATIONS
%----------------------------------------------------------------------------------------

\documentclass[12pt]{article}
\usepackage[english]{babel}
\usepackage[utf8x]{inputenc}
\usepackage{amsmath}
\usepackage{graphicx}
\usepackage[colorinlistoftodos]{todonotes}
\usepackage{ragged2e}
\usepackage[none]{hyphenat}
\usepackage[hidelinks]{hyperref}
\usepackage{listings}
\usepackage[title,titletoc,toc]{appendix}
\usepackage{kantlipsum}
\usepackage{spverbatim}
\usepackage{caption3}
\usepackage{longtable}
% This is for adding extra subsections
\usepackage{titlesec}
\titleclass{\subsubsubsection}{straight}[\subsection]
\newcounter{subsubsubsection}[subsubsection]
\renewcommand\thesubsubsubsection{\thesubsubsection.\arabic{subsubsubsection}}
\renewcommand\theparagraph{\thesubsubsubsection.\arabic{paragraph}}
\renewcommand\thesubparagraph{\theparagraph.\arabic{subparagraph}}

\titleformat{\subsubsubsection}
  {\normalfont\normalsize\bfseries}{\thesubsubsubsection}{1em}{}
\titlespacing*{\subsubsubsection}
{0pt}{3.25ex plus 1ex minus .2ex}{1.5ex plus .2ex}

\makeatletter
\renewcommand\paragraph{\@startsection{paragraph}{5}{\z@}%
  {3.25ex \@plus1ex \@minus.2ex}%
  {-1em}%
  {\normalfont\normalsize\bfseries}}
\renewcommand\subparagraph{\@startsection{subparagraph}{6}{\parindent}
  {3.25ex \@plus1ex \@minus .2ex}%
  {-1em}%
  {\normalfont\normalsize\bfseries}}
\def\toclevel@subsubsubsection{4}
\def\toclevel@paragraph{5}
\def\toclevel@paragraph{6}
\def\l@subsubsubsection{\@dottedtocline{4}{7em}{4em}}
\def\l@paragraph{\@dottedtocline{5}{10em}{5em}}
\def\l@subparagraph{\@dottedtocline{6}{14em}{6em}}
\@addtoreset{subsubsubsection}{section}
\@addtoreset{subsubsubsection}{subsection}
\@addtoreset{paragraph}{subsubsubsection}
\makeatother
\setcounter{secnumdepth}{6}
\setcounter{tocdepth}{6}
%-----------------------------------------------%

\usepackage{array}
\newcolumntype{L}[1]{>{\raggedright\let\newline\\\arraybackslash\hspace{0pt}}m{#1}}
\newcolumntype{C}[1]{>{\centering\let\newline\\\arraybackslash\hspace{0pt}}m{#1}}
\newcolumntype{R}[1]{>{\raggedleft\let\newline\\\arraybackslash\hspace{0pt}}m{#1}}

\begin{document}

\begin{titlepage}

\newcommand{\HRule}{\rule{\linewidth}{0.5mm}} 
\center

%----------------------------------------------------------------------------------------
%	LOGO SECTION
%----------------------------------------------------------------------------------------

\includegraphics{images/uva.jpeg}\\[0.5cm]% Include a department/university logo - this will require the graphicx package
 
%----------------------------------------------------------------------------------------

%----------------------------------------------------------------------------------------
%	HEADING SECTIONS
%----------------------------------------------------------------------------------------
\textsc{\Large System and Network Engineering, MSc}\\[0.5cm] 
\textsc { \large Research Project 1}\\[0.4cm] % Title of your document

%----------------------------------------------------------------------------------------
%	TITLE SECTION
%----------------------------------------------------------------------------------------
\HRule \\[0.4cm]
{ \huge \bfseries Security Intelligence Data Mining}\\[0.4cm] % Title of your document
{ \large \bfseries Research Proposal}\\[0.4cm] % Title of your document
\HRule \\[0.4cm]




%----------------------------------------------------------------------------------------
%	AUTHOR SECTION
%----------------------------------------------------------------------------------------


\large Diana Rusu\\
{\bfseries Diana.Rusu@os3.nl}\\[0.5cm]
\large Nikolaos Petros Triantafyllidis\\
\bfseries Nikolaos.Triantafyllidis@os3.nl\\[2cm]

{\large \today} 

\end{titlepage}

\newpage
\section*{Abstract}
\newpage
\section*{Acknowledgement}

We would like to express our gratitude and appreciation for all support, expert knowledge and guidelines during this research project. It has been a great opportunity for us to explore the real feel of working with a successful company. 
\\
\\
Who should we mention? 
\begin{itemize}
\item{dhr. prof. dr. ir. C.T.A.M. (Cees) de Laat - for the project proposals}
\item{Delloite NL specifically to our first supervisor Henri Hambartsumyan - for giving us access inside their company and offices, for all the experts and employees that we had the chance to meet}
 \begin{itemize}
 \item{Henri Hambartsumyan - who proposed Security Intelligence Data Mining}
 \item{Kremers, Joost  - our main supervisor for guideline during the research }
 \item{Niels - Cyber Threat Intelligence project}
 \item{Ari Davies - proposed making the plugins also more depth in CTI }
 \item{Gijs Hollestelle - presented the actual CTI platform and the portal to it, proposed phishtank and cleared out the purpose for this project(making some agents for IR system) }
 \end{itemize}
\item{dhr. dr. C.P.J. (Karst) Koymans}
\end{itemize}
\newpage

\tableofcontents
\newpage

\section{Introduction}
%\addcontentsline{toc}{section}{Introduction}
\parbox{\linewidth}{
With the increasing number of cyber-attacks and the growth of computer crime worldwide, it becomes apparent that IT security is a major concern and crucial survival factor for large companies, organisations and institutions of any sort. Security Operations departments working to ensure confidentiality, integrity and availability for the system infrastructure of their organisation, invest huge  parts \cite{cyber} of their time and effort in detecting threats in real time. A very valuable source of security intelligence, vital to cyber-risk assessment, is information mined from data posted on public sites such as "pastebins" or social networks. However, this is a very cumbersome task due to the lack of Natural Language Processing capabilities in most of the existing tools. Moreover, as recent events have showcased \cite{sony}, several threats arise from governments and criminal associations originating from countries whose languages use non-latin scripts (Chinese, Russian, Korean, etc.). It is, hence, important to have data mining tools that provide support for such alphabets and languages, since a lot of a security intelligence can be discovered in such texts. The main goals of this research project will be to explore the various public data and detect the most appropriate among them. Moreover, numerous current data analytics techniques as well as their application on security related issues will be assessed. Lastly the above knowledge will be applied on the implementation of a simple system that will work as a proof-of-concept and help determine the technical feasibility, storage requirements and operational cost of such a system. This project was proposed by and will carried out in co-operation with Deloitte NL. 
}
\newpage
\subsection{Research Questions}
This topic is admittedly very open but it can be narrowed down to several specific research questions some of which we will try to answer to some extent. The main question that we will be trying to answer is the following:\\[0.1cm]

\noindent
\textbf{How can we effectively use public sources to obtain real time information about security incidents?}\\[0.1cm]

\noindent
This question can be analysed into more specific parts that cover the topic to some extent, as follows:

\begin{enumerate}
	\item How can the raw data be effectively collected from the public sources? 
	\begin{itemize}
		\item How can we effectively detect the reliable sources?
		\item What search terms can we deploy during the retrieval phase?
		\item How can the unstructured data be pre-processed? 
	\end{itemize}
	\item How can the data be analysed in respect to security operations?
	\begin{itemize}
		\item How can we apply current Data Mining and Analytics techniques on Security issues?
		\item How can we derive the risk assessment model from the above?
		\item How can we apply the model on new data?
	\end{itemize}
	\item	How can the collected knowledge be applied on a system implementation?
	\begin{itemize}
		\item What is a reliable and extensible System Architecture that can be designed?
		\item What are the computational and storage requirements of such a system?
		\item What extensions can be proposed for that system?
	\end{itemize}
\end{enumerate}

The proposed system extensions can spawn further research questions, namely on the topics of presenting the analysed data, reacting to the real time events and finally assessing the situations that arise and providing feedback to the system.


\subsection{Related work}
There is a lot of literature around the field of Data Mining and more recently Web Mining. The most prominent and recent case is the book 'Mining The Social Web' by M.Russel \cite{socialweb} that deals with exploring and mining information from social websites (e.g., Facebook, Twitter, LinkedIn, Google+, GitHub, etc.). There are also several academic papers and books that deal with applying Data Mining to System Security. One example is a system proposed by the university of Minnesota, called MINDS, that employs various Data Mining in Intrusion Detection. The system is described in their paper 'Data Mining for Cyber Security' \cite{minds}.  Another example is a system proposed by the Dutch company Sentient in co-operation with the Amsterdam Police Force \cite{police} aiming to provide Data Analytics operations automation while on the same time minimising the technical expertise needed by the system user. 
\subsection{Ethical implications}
%\addcontentsline{toc}{section}{Ethical implications}
The main part of this research comprises of exploring current techniques and their application on IT security, as well as the specification of a system that employs Data Mining techniques to collect security intelligence. In order for the models to be defined some amount of information will have to be gathered. This information will originate solely from public sources and will be mostly historical data. In the unlikely case that any previously unnoticed security issues are encountered they will be handled with discretion and communicated only towards the appropriate targets and only with the approval of the OS3 core team and Deloitte Digital. The collection or storage of personal data is not intended and any collected information will be discarded after the end of this project. The usage of shared computational and network infrastructure will only be used for the needs of this project and within the legal limits.

\section{Methodology}
WRITE SMALL INTRO HERE
This project was carried out in collaboration with Deloitte Risk Services, Netherlands. The first step is to get acquainted with the Cyber-Risk Services and Data Analytics teams of the company and their people. When directed to the appropriate experts on each respective field (Security, Data Analytics, NLP, etc.) we will be working closely with them to collect the needed knowledge and expertise in order to reduce an effective risk assessment model. This will be made possible by applying Data Mining techniques that are already in use in other fields (economic risk assessment, marketing, etc.) on security issues. In parallel we will start designing and implementing the Information Retrieval part of the system. Once the appropriate amount of Data has been gathered from public sources we will apply the appropriate techniques that have been proposed to us on the collected dataset in order to design the risk assessment model. The next step will be to apply the designed model on new data and evaluate the effectiveness of the model. Lastly, based on the gathered experience and knowledge we will try to design a modular and extensible system architecture that deploys the designed techniques. Based on that design we will attempt to implement a simple proof of concept that will help us define the resource requirements and operational cost of such a system. 

\subsection{Hardware \& Software Used}
\parbox{\linewidth}{
Most of the work was carried out on end laptops using the Guest WiFi network of the Deloitte headquarters for internet connectivity. For the execution of the demonstration software and storage of the collected data, one of our assigned OS3 servers was used (oxford.studlab.os3.nl).\\ 
For the storage of the raw well as the processed data we used MongoDB, a free and open source NoSQL database system.\\ 
All software developed for the purposes of this project was written in the Python programming language using standard as well as external libraries. The external libraries and tools used were Tweepy \cite{tweepy}, Scrapy \cite{scrap}, Numpy \cite{numpy}, PyMongo \cite{pymongo}, Orange \cite{orange}, scikit-learn \cite{sclera}, NLP-Toolkit \cite{nltk}. More details on their usage can be found in the next chapters, where each software module is explained.}

\subsection{Proposed System}
As an outcome of our research we propose a system that will be able to aggregate information from the public web which it will preprocess and analyse in order to gather Intelligence regarding IT Security incidents and raise real time alerts. 
 
\subsubsection{General Outline}
The proposed system will be modular in the sense that there will be standalone software processes responsible for a very specific task. The output of each module feeds the next module in succession thus forming a pipeline of modules that compose the whole system. \\
Moreover we propose the integration of the system under a central Control Portal that will be responsible for the initial configuration and overview of the status of each Software module.\\
We distinguish five different steps in the system each one of them corresponding to a different software module. These are:

\begin{enumerate}
	\item \textbf{Configuration}\\ 
	The System Operator defines a set of configuration options (keywords, data sources, thresholds etc.) that will be propagated to each separate module.
	\item \textbf{Data Aggregation}\\
	The agents belonging to the Aggregator module collect raw data from their designated sources. At this point a first structure is given to the raw data collected before they are stored to the corresponding database. 
	\item \textbf{Data Deflation \& Filtering}\\
	A common format is given to all the data collected from different sources. We have to be able to distinguish the interesting from the uninteresting documents so at this point some sort of scoring system is applied for each document collected. The documents that have been distinguished are then inserted into the warehouse in a format suitable for further processing.
	\item \textbf{Analytics \& Alerting}\\
	A series of Data Mining operations is run against the collected dataset to extract previously unnoticed patterns. Based on the outcome, and after the application of certain predefined criteria, the corresponding module will have to be able to raise alerts and notify the appropriate parties for any probable security incident in progress.
	\item \textbf{Feedback \& Reconfiguration}\\
	A Security Analyst checks the raised alerts and assesses the real risk. The divergence from the estimated risk is calculated and according to that a set of readjustments is propagated to the configurable parts of the system.
\end{enumerate}

\noindent Having introduced the last step it becomes apparent that the system operates as a feedback loop since the system returns to the configuration step after assessment, in order to achieve better accuracy. \\
Each part of the system is described in detail in the following sections. All steps taken towards building an initial proof of concept for this system are described as well. 

\subsubsection{Configuration}
\subsubsubsection{Specification}
\subsubsubsection{Implementation}
\subsubsection{Data Aggregation Module}
This chapter describes a system for collecting and mining information from public sources. First subsection will define the relevant websites that provide interest in detecting threats, malwares and phishing sites. For each specific source, next subsections will investigate API functionality and steps in implementing crawling part.
\subsubsubsection{Specification}

Discovering real time threats with a plausible margin of error, became a challenging task as Internet provides an enormous database with random messages and information. Among these sites, significant information related to Security Intelligence can be found. Therefore it is important to define reliable sources that being monitored can provide significant valuable data for detecting future or present possible attacks. Special attention has to be provided in establishing which are the optimum sources that will provide the seeking information. 

What sources can be considered as reliable? In order to answer this question inside research has been given by experts. From a company perspective it is important to have an overview over all general possible attacks or client specific. Among all the possible feed sources, four are considered at this moment mostly valuable. Other sources that provide RSS feeds or IRC channels will not be investigated due to time limitations and company requirements. This research will primarily focus on the following sites: Phishtank, Pastebin, Twitter and Reddit. Mainly, these are considered fundamental for the reason that past activity revealed vulnerabilities found through posted messages.  Furthermore, mentioned sources will be analysed in particular and be judged by the relevance of the subjects involved in discussion. 
\subsubsubsection{Implementation}
\paragraph*{Twitter}
\addcontentsline{toc}{paragraph}{Twitter}
\paragraph*{Pastebin}
\addcontentsline{toc}{paragraph}{Pastebin}
\hfill \break
"Pastebins" websites \cite{fpaste} \cite{pastebin} allow everyone with or without registration to share real time text or code snippets. It attracted many users during past years including cybercriminals like malware writers\cite{pastebin-magazine}. The enormous flow of information include from database dumps containing e-mails and passwords to harmful backdoor programs. At a deeper examination of these public pasted messages possible future attacks can be determined. From all possible "pastebins" sources (e.g fpaste.org, paste2.org, pastie.org\cite{pastebin-pastie}, paste.ubuntu.org.cn \cite{pastebin-ubuntu} etc.) pastebin.com \cite{pastebin} will be further inspected and monitored.

Pastebin API offers different options for developers, more specifically, creating new pastes, listing trending pastes and pastes from a particular user. This extensions can be used once one obtains the unique Developer API Key. 

At this moment Pastebin API \cite{pastebin} does not offer complete information retrieval option for listing all new messages that arrive from unregistered users. The option available would be to list messages from registered users, though all information provides interest. This makes the crawling module challenging but not impossible. A method will be introduced to overcome this limitation. 

Search method for retrieving information using keywords uses third party software namely, "Google Custom Search". Therefore, if one desires to obtain full database it would have to use Google's API. Free users account will be limited to 1000 pastes per day. 

To overpass API extension limitation and mine all new pastes that arrive in real time a web crawler module will be implemented. Python offers an open source framework called "scrapy" \cite{scrapy} for building an extensible webcrawler. Pastebin's \texttt{Archive} page contains latest pastes posted in the last 10 minutes at least. We will crawl and fetch this page each 10 minutes and store data into mongoDB database. Not all data from the page provides interest this is why we will scrape only information that is relevant for further examination. This includes the url, paste, the date it was posted and the number of unique views. By inspecting the elements required from the DOM tree, we then can parse this nodes and extract only the text. Once all this data has been gathered, it will be dumped in the database.
\paragraph*{Phishtank}
\paragraph*{Reddit}
\addcontentsline{toc}{paragraph}{Reddit}
\hfill \break
Reddit community comprises social websites and news sites. It contains over 5000 channels called "subreddits" which belong to different categories. This research will mine Reddit \cite{reddit} data from the following branches:
\begin{itemize}
\item /r/blackhat \cite{red-blackhat}
\item r/malware \cite{r.malware}
\item r/netsec \cite{r.netsec}
\item r/pwned \cite{r.pwned}
\item /r/vrd/ \cite{r.rvd}

\end{itemize}

Reddit API\cite{reddit} provides a rich list of methods which are easy to use. For listing new, top, controversial etc. messages is sufficient to add "/" to the subreddit one wants to obtain information. Endpoints can be used to filter and sort posts namely, \textit{after/before, limit, count} and \textit{show} \cite{reddit}. This parameters will be later used in retrieving lists of messages in json format. 

All ingredients for implementing a retrieval agent for Reddit are supplied by its API. The only part we need to do is to filter and obtain data that further investigated and processed can help defining the relevance of this information. First a connection needs to be opened to the the targeted subreddit, in this case to /r/blackhat. 
\begin{spverbatim}
response = urllib2.urlopen('http://www.reddit.com/r/blackhat/new
			.json?sort=new&limit=100')
\end{spverbatim}
\hfill \break
Response variable will retrieve a file in json format, from which the title associated to the url, number of comments and score will be stored, because this information provides interest for further processing. Python provides an open json library for processing this type of files. All scraped Reddit data will be stored in mongoDB database. 

\paragraph*{Trigger System}
\addcontentsline{toc}{paragraph}{Trigger System}
\subsubsection{Deflation/Warehousing/Filtering}
\subsubsubsection{Specification}
\subsubsubsection{Implementation}
\subsubsubsection{Alternatives}
A strict method for deflating documents has been previously introduced. This subsection will outline as alternatives known processes for cleaning documents by removing stop words, stemming and recognizing part of speech with a further extension to TF-IDF alghorithm for determining the relevance and weights for the documents.
\paragraph*{NLTK}
NLTK or Natural Language Toolkit ...

\paragraph*{TF-IFD}
\hfill \break
\\
TF-IDF \cite{tf-idf} stands for term frequency–inverse document frequency and it is used to statistically score the importance of a word in a document or across a collection of documents. A word will be considered important if it appears multiple times into a document. TF is calculated by counting how many times the specific word appears, over the total number of words from that document and it will get a high score. But if it appears and in other documents it will not be as unique and will receive a lower score. 
\\
\\
%Having this in mind we implemented and applied the algorithm  on the retrieved data but the results turned out to be not sufficient to grade the value of the documents.
For determining the most relevant documents, by defining query terms which are related to the known threats categories, these will be further processed by the algorithm and as a result a statistical list of the most significant documents will be obtained. This is an automated method for sorting document that are more appropriate to a specific queried subject, given on-topic keywords list.

\subsubsection{Analytics and Mining}
This section will sketch few methods that can be applied over the congregated information. All the unstructured data has to be processed and converted to "clean" structured information which includes natural language processing. For this step Python provides an open source library nltk \cite{nltk}. We have to dip into several phases to get a good transformation of  paste, mesages or titles. This includes stop words removal, stemming and pos(part of speach) detection. 
At this moment there are various data mining techniques \cite{oracle-list} that include text processing. Investigating all this techniques is out of scope in this project, therefore the once that are mostly commonly used will be provided. In data mining step or knowledge discovery from database(KDD) it is important to find and extract patterns and knowledge from database not the data itself \cite{data-kdd}.
\subsubsubsection{Specification}
\subsubsubsection{Implementation}
\subsubsection{Feedback and Assessment}
\subsubsubsection{Specification}
\paragraph{Data mining tasks}
\begin{itemize}
\item 
Classification: Refers to the task of generalising known structures and applying them to new
data. For example, an e-mail program might attempt to classify an e-mail as "legitimate" or as
“spam”. Since this involves training of the algorithm before use, it might not be ideal for use in
threat discovery. However, it can be very useful for providing feedback to the system and fine
tuning it to be able to classify the documents for each threat category, acting together with the
configuration step
\item 
Anomaly detection (Outlier/change/deviation detection): This task refers to identification of
unusual data records that might imply unusual behaviour in the dataset. For example, in our
use case, a first indication that something might be possibly wrong would be an exceptional
mention of a company name within a certain period of time. The problem here is to determine
what can be defined as exceptional. This of course is not possible without the use of old data
so a certain threshold can be determined. Until sufficient amount of data has been gathered we
can naively implement this by setting the thresholds manually. This operation is particularly
popular in intrusion detection systems.
\item 
Association rule learning (Dependency modelling): This task refers to the search for
relationships between variables. One example of this is the infamous beer and diapers story.
This test requires no training set and can be performed on any structured set of data. In our
use case this can help us determine trends between the occurrences of keywords, such as the
name of a company together with the name of a name of an attack. The output of this test also
contains metrics about the support (% of the documents supporting the association) as well as
confidence (ratio of the support percentage of either part of each rule).
\item 
Clustering: Refers to the task of discovering groups and structures in the data that present
some sort of similarity, without using known structures in the data. In our case this can help us
determine new categories of documents that have not been predefined (remember the threat
categorisation we mentioned in the beginning).
\item
Regression: This task attempts to find a function which models the data with the least error.
This might not be directly applicable to our specific example but based on the experience
gathered after a certain period of time that the system will be operated, we can attempt to
design an analytic model using regression analysis techniques.
\item
Summarisation: This task aims to provide a more compact representation of the data set,
including visualisation and report generation. In our case this applies to the creation of the
alerts to be presented to the Security Analysts. The alerting will consist of a common format
that will apply to automatically generated events that will be propagated to the appropriate
parties. One proposed format could be the following:
\{alert\_id: 00001, subject: "Something is rotten in the state of Denmark", importance: Red,
%backing_documents:\[1,3,4,6\]\}
This presents the analyst with an event informing them about a security incident with a certain
importance level. The way of determining the level of importance falls out of the scope of this
document. Notice also the list of supporting documents resulting from different sources. At this
point the analyst can go and inspect each one of this documents manually, assess the risk
themselves and take the appropriate measures such as further informing the involved parties, or
ignoring the alert.
\item Attribute Important 
\item etc.
\end{itemize}
\subsubsubsection{Implementation}
\paragraph*{Association Rule}
\addcontentsline{toc}{paragraph}{Association Rule}
\paragraph*{Clustering}
\addcontentsline{toc}{paragraph}{Clustering}
\hfill \break 
\\
For clustering N collected documents from the database in K categories an option would be, by using K-means\cite{k-means} algorithm also known as Lloyd's algorithm. The method\cite{k-means-example} proposed by this algorithm aims to find K non-overlapping clusters and it works as follows. User specifies how many K clusters will represent the initial centroids, in which the information will be divided into. Assign each observation point/word to the closest centroid, where each collection of points/words form a cluster. The centroid is then updated depending of the words added to this cluster. The process stops when the words assigned will not change the clusters. At this moment there is no guarantee that the result retrieved is optimum. The outcome might be different each time, depending on how the initial centroids are chosen. Important factor to mention here is the selection of K, because an incorrect choice might lead as well to non-optimum results. There are several procedures\cite{procedures-for-kmeans} to calculate a favourable number of clusters but this is out of scope and the investigation will consider random cluster number that will be feasible chosen compared to the amount of data.
\\
\\
In order to implement the proposed algorithm an open  source scikit- learn \cite{sklearn} tool was used for data analysing that provides a library called SKlearn which can be integrated in Python environment. Following an example provided in a tutorial from the sklearn website, which attemps to categorize 20 news group\cite{k-means-20news}, we adapted it to use the data provided in the database and show as well what are the documents that are related to each other in a cluster.  In this case we applied the algorithm on over 600000 twitter messages and the expected outcome was to categories this messages in 20 groups. From the clusters obtained two showed big interest as it presents real events that happened in the past weeks.
\begin{spverbatim}
Cluster 13:  paris  charlie  hebdo  attack  http  mayor  nypd  
rt  victims  french  honor  visited  nyc  france  terror  muslim 
 bolsters  security  jewish  cover
Cluster 15:  photos  leaked  upton  kate  jennifer  lawrence 
 nude  victoria  justice  megan  fox  http  seen  rt  hoeuu2dubr  
itweetlikegirls  hilarious  kardashian  kim  aigbtbpmvv
\end{spverbatim}
\hfill \break
Observing the resulted clusters we can claim that the algorithm behaved well considering that it is an unsupervised method and the top frequent terms describe in Cluster 13 the most recent attack in Paris, while the Cluster 15 which is totally unrelated to this attack outlines the incident with the leaked photos of celebrities.
\paragraph*{Classification}
\section{Results}
\subsection{Demo}
\section{Conclusions}

\section{Future Work}

At this moment we have a system 

\newpage
\section{Appendix}
\newpage

\begin{thebibliography}{99}
\bibitem{cyber}
US cybercrime: Rising risks, reduced readiness, PWC, Available at: \url{http://www.pwc.com/en_US/us/increasing-it-effectiveness/publications/assets/2014-us-state-of-cybercrime.pdf}
\bibitem{sony}
	 D. Sanger, N. Perlroth.  U.S. Said to Find North Korea Ordered Cyberattack on Sony. [online] Nytimes.com. \\Available at: \url{http://www.nytimes.com/2014/12/18/world/asia/us-links-north-korea-to-sony-hacking.html?_r=0}
\bibitem{socialweb}
  Russell, MA (2014). Mining the Social Web, O'Reily Media, USA
\bibitem{minds}
    V. Chandola et al. Data Mining for Cyber Security, Department of Computer Science, University of Minnesota, Springer, 2006
\bibitem{police}
   RCP van der Veer, H.T. Roo,  A. van der Zanden, Data mining for intelligence led policing, Sentient, Amsterdam Police Force, Amsterdam, The Netherlands, 2009
\bibitem{nltk}
Nltk.org, (2015). Natural Language Toolkit - NLTK 3.0 documentation. [online] Available at: \url{http://www.nltk.org/}

\bibitem{fpaste}
Fedora Project Pastebin Fpaste.org, (2015). New paste • Fedora Project Pastebin. [online] Available at:  
\url{http://fpaste.org }
\bibitem{pastebin}
Pastebin, (2015). Pastebin.com - \#1 paste tool since 2002!. [online] Available at:     \url{http://pastebin.com/}

\bibitem{pastebin-magazine}
Security Intelligence, (2015). Pastebin a Convenient Way for Cybercriminals to Remotely Host Malware. [online] Available at: \url{http://securityintelligence.com/news/pastebin-convenient-way-cybercriminals-remotely-host-malware/#.VMS25DX8vCI }
\bibitem{pastebin-pastie}
Pastie.org, (2015). New - Pastie. [online] Available at: \url{http://pastie.org} 
\bibitem{pastebin-ubuntu}
 Paste.ubuntu.org.cn, (2015). [online] Available at: \url{http://paste.ubuntu.org.cn/}
\bibitem{scrapy}
Scrapy.org, (2015). Scrapy | A Fast and Powerful Scraping and Web Crawling Framework. [online] Available at: \url{http://scrapy.org/} 
\bibitem{reddit}
Reddit.com, (2015). reddit: the front page of the internet. [online] Available at: \url{http://www.reddit.com/} 
\bibitem{red-blackhat}
Reddit, (2009). blackhat library: Documenting blackhat hacking techniques • /r/blackhat. [online] Available at: \url{https://www.reddit.com/r/blackhat}
\bibitem{r.malware}
Reddit, (2009). Malware Analysis \& Reports • /r/Malware. [online] Available at:   \url{http://www.reddit.com/r/malware}
\bibitem{r.netsec}
reddit, (2007). /r/netsec - Information Security News \& Discussion. [online] Available at: \url{https://www.reddit.com/r/netsec/} 
\bibitem{r.pwned}
Reddit, (2008). pwned • /r/pwned. [online] Available at: http://www.reddit.com/r/pwned 
\bibitem{r.rvd}
Reddit, (2012). Vulnerability Research and Development • /r/vrd. [online] Available at: \url{http://www.reddit.com/r/vrd/}
\bibitem{oracle-list}
Oracle.com, (2015). Oracle Data Mining Techniques and Algorithms. [online] Available at: \url{http://www.oracle.com/technetwork/database/enterprise-edition/odm-techniques-algorithms-097163.html}
\bibitem{tf-idf}
Tfidf.com, (2015). Tf-idf : A Single-Page Tutorial - Information Retrieval and Text Mining. [online] Available at: \url{http://www.tfidf.com/} 
\bibitem{data-kdd}
Han, Jiawei; Kamber, Micheline (2001). Data mining: concepts and techniques. Morgan Kaufmann. p. 5. ISBN 9781558604896. "Thus, data mining should have been more appropriately named "knowledge mining from data," which is unfortunately somewhat long"
\bibitem{owasp}
Owasp.org, (2015). Category:OWASP Top Ten Project - OWASP. [online] Available at: \url{https://www.owasp.org/index.php/Top10#OWASP_Top_10_for_2013} 
\bibitem{list-2015-attacks}
Passeri, P., Passeri, P. and Passeri, P. (2015). Cyber Attacks Timeline | Hackmageddon.com. [online] Hackmageddon.com. Available at: http://hackmageddon.com/category/security/cyber-attacks-timeline/ 
\bibitem{sklearn}
Scikit-learn.org, (2015). scikit-learn: machine learning in Python — scikit-learn 0.15.2 documentation. [online] Available at: \url{http://scikit-learn.org/stable/index.html}
\bibitem{orange}
Bioinformatics Laboratory, U. (2015). Orange Data Mining. [online] Orange.biolab.si. Available at: \url{http://orange.biolab.si/} 
\bibitem{k-means} Junjie Wu , Advances in K-means Clustering: A Data Mining Thinking
\bibitem{k-means-example}
Gonçalves, H. (2015). K-means clustering - algorithm and examples. [online] Onmyphd.com. Available at: \url{http://www.onmyphd.com/?p=k-means.clustering}
\bibitem{procedures-for-kmeans}    Glenn W. Milligan, Martha C. Cooper , An examination of procedures for determining the number of clusters in a data set
\bibitem{k-means-20news}
Scikit-learn.org, (2015). Clustering text documents using k-means — scikit-learn 0.15.2 documentation. [online] Available at: \url{http://scikit-learn.org/stable/auto_examples/document_clustering.html} 
 
\end{thebibliography}
\end{document}